\documentclass{situationist-preprint}

\title{Leibniz's Relationalism Reconsidered}
\msid{TSM-06-b}
\msfamily{TSM-06}
\msdiscipline{Time, Space and Modality}
\mstype{Companion manuscript}
\mskeywords{Leibniz; relationalism; Leibniz--Clarke correspondence; force; absolute motion; substantivalism}

\begin{document}
\maketitle

\begin{abstract}
Leibniz's arguments against absolute space are the classical case for relationalism, and TSM-06-a answers them briefly. We reconsider his position. His arguments from sufficient reason and the identity of indiscernibles succeed against Newton's absolute space, but only against its absolute positions and velocities, which spacetime substantivalism discards. His own account of true motion, as motion whose cause is in the body, either fails to distinguish rotation from its relational duplicate or presupposes the affine structure that substantivalism asserts. Leibniz is right about rest and wrong about space.
\end{abstract}

\section{What Leibniz's arguments show}

In the correspondence with Clarke, Leibniz argues that if space were a substance, God could have placed the whole universe elsewhere, or set it in uniform motion, without any discernible difference; such a choice would lack a sufficient reason, and the two arrangements would be indiscernible and so identical \citep{Leibniz1716}. The arguments are sound against what they target. Newtonian absolute space contains facts about absolute position and absolute velocity that make no difference to anything observable.

They are not sound against substantivalism about spacetime structure (Definition~2.1 of TSM-06-a). Neo-Newtonian spacetime has an affine structure, which distinguishes accelerated from unaccelerated motion, and no absolute rest \citep{Stein1967,Earman1989}. A static shift or a uniform boost of the whole universe is not a difference in it. Leibniz's arguments thus remove the structure the substantivalist no longer posits, and leave the structure the bucket argument requires.

\section{Leibniz's account of true motion}

Leibniz conceded that true motion differs from mere change of relative situation, and located the difference in the cause of motion: a body truly moves when the cause of the change is in it \citep[fifth letter, §53]{Leibniz1716}. The concession is the crux of his position.

\begin{proposition}\label{p:dilemma}
Either the force that marks true motion is fixed by the relations among bodies, in which case it does not distinguish the rotating bucket from its relational duplicate, or it is not, in which case it requires structure beyond those relations of the kind that Lemma~3.2 of TSM-06-a identifies.
\end{proposition}

\begin{proof}
Consider the rotating bucket and a duplicate in which the bucket is at rest and the rest of the universe rotates around it. The relations among bodies are the same in both, and the dynamics distinguish them by the concavity of the water (Principle~P1 of TSM-06-a). If force is fixed by the relations among bodies, it is the same in both, and it does not mark the difference that Leibniz admitted. If it is not, then the forces in the two cases differ although the relations do not, and their difference is a difference in which bodies deviate from unaccelerated motion. That a body deviates from unaccelerated motion is a fact about its trajectory relative to the affine structure, which by Lemma~3.1 of TSM-06-a relative motion cannot supply.
\end{proof}

Leibniz's dynamics treated force, not motion, as real, and this is sometimes read as a primitive-force relationalism \citep{Roberts2003}, of the kind that \citet{Sklar1974} later proposed as primitive acceleration. TSM-06-a answers Sklar (Objection~3 there): primitive accelerations must cohere across bodies and their parts, as if all unaccelerated bodies moved along one common family of straight lines, and this coherence is explained by a single affine structure and left a coincidence without it. The same reply applies to Leibniz's force. His concession, taken seriously, is a concession to substantivalism about the affine structure of spacetime.

\section{Conclusion}

Leibniz's arguments from sufficient reason and indiscernibility are correct against absolute rest and inapplicable to spacetime structure. His account of true motion either fails to explain the bucket or presupposes affine structure. Reconsidered, his position supports the conclusion of TSM-06: Newton's bucket survives, and space, as spacetime structure, is a substance.

\begin{thebibliography}{9}
\bibitem[Earman(1989)]{Earman1989} Earman, J. (1989). \emph{World Enough and Space-Time: Absolute versus Relational Theories of Space and Time}. MIT Press.
\bibitem[Leibniz and Clarke(1716)]{Leibniz1716} Leibniz, G.~W., and Clarke, S. (1715--1716). \emph{The Leibniz--Clarke Correspondence}. Ed. H.~G. Alexander, Manchester University Press, 1956.
\bibitem[Roberts(2003)]{Roberts2003} Roberts, J.~T. (2003). Leibniz on force and absolute motion. \emph{Philosophy of Science}, 70(3), 553--573.
\bibitem[Sklar(1974)]{Sklar1974} Sklar, L. (1974). \emph{Space, Time, and Spacetime}. University of California Press.
\bibitem[Stein(1967)]{Stein1967} Stein, H. (1967). Newtonian space-time. \emph{The Texas Quarterly}, 10, 174--200.
\end{thebibliography}

\end{document}
